\section{Praktische Aspekte digitaler Regler}

\subsection{Pulsweitenmodulation (PWM)}
\begin{tcolorbox}
\boxtitle{Funktionsprinzip}
Zur verlustarmen Leistungssteuerung wird die feste Versorgungsspannung $U_0$ mit konstanter Schaltfrequenz $f_{\text{PWM}} = \frac{1}{T_{\text{PWM}}}$ getaktet.\\[2pt]
\textbf{Tastgrad (Duty Cycle):}
\[
v = \frac{T_{\text{ein}}}{T_{\text{PWM}}} \in [0, 1] \quad (\text{bzw. } 0 \dots 100\%)
\]
\textbf{Mittelwertbildung} (Trägheit von Motor oder Last $\gg T_{\text{PWM}}$):
\begin{align*}
\text{Mittlere Spannung:} \quad &\bar{u} = v \cdot U_0 \\
\text{Mittlerer Strom:} \quad &\bar{i} = v \cdot I_0 \\
\text{Mittleres Drehmoment:} \quad &\bar{M} = v \cdot M_0
\end{align*}
\textbf{Vorteil:} Schalter (MOSFETs) sind voll durchgeschaltet oder gesperrt $\implies$ minimale Verlustleistung im Stellglied!
\end{tcolorbox}

\begin{tcolorbox}
\boxtitle{H-Brücken-Ansteuerung (4-Quadrantenbetrieb)}
Ermöglicht Drehrichtungsumkehr und aktives Bremsen:
\begin{center}
\small
\begin{tabularx}{\columnwidth}{l c c c c c}
\toprule
\textbf{Zustand} & \textbf{DIR} & \textbf{PWM} & \textbf{ENA} & \textbf{OUT1} & \textbf{OUT2} \\
\midrule
Antrieb pos. (Puls) & L & H & H & L & H \\
Antrieb neg. (Puls) & H & H & H & H & L \\
Bremse low-side (Pause) & L & L & H & L & L \\
Bremse high-side (Pause) & H & L & H & H & H \\
Freilauf / Abschaltung & X & X & L & Z & Z \\
\bottomrule
\end{tabularx}
\end{center}
\scriptsize ($H$: High, $L$: Low, $X$: beliebig, $Z$: hochohmig/gesperrt).
\end{tcolorbox}

\subsection{Integer- und Festkomma-Arithmetik}
\begin{tcolorbox}
\boxtitle{Gründe für Festkomma-Arithmetik}
Einfache Mikrocontroller besitzen keine Gleitkomma-Hardware (FPU). Software-Gleitkomma erhöht Codeumfang und Rechenzeit drastisch.
\end{tcolorbox}

\begin{tcolorbox}
\boxtitle{Festkomma-Skalierung (Q-Format)}
Eine Fließkommazahl $x$ wird durch Multiplikation mit $2^Q$ als Ganzzahl dargestellt:
\[
x_{\text{int}} = \text{round}(x \cdot 2^Q)
\]
\begin{itemize}
  \item \textbf{Addition / Subtraktion:} Direkt möglich (gleiches $Q$-Format vorausgesetzt).
  \item \textbf{Multiplikation:}
  \[
  z_{\text{int}} = (x_{\text{int}} \cdot y_{\text{int}}) \gg Q \quad (\text{Rechtsshift um } Q \text{ Bits})
  \]
  \item \textbf{Zwischenergebnisse:} Immer in doppelter Bitbreite berechnen (z.\,B. \texttt{int16\_t} $\times$ \texttt{int16\_t} $\to$ \texttt{int32\_t}).
\end{itemize}
\end{tcolorbox}

\begin{tcolorbox}
\boxtitle{Typische numerische Gefahren}
\begin{enumerate}
  \item \textbf{Zahlenbereichsüberlauf (Overflow):}
  Führt bei Zweierkomplement zu Vorzeichenumkehr (z.\,B. $+127 + 1 = -128$ bei \texttt{int8}). $\implies$ \textbf{Sättigungsarithmetik} einsetzen!
  \item \textbf{Kürzungsfehler bei Division:}
  Eine Division vor einer Multiplikation kann ganzzahlig zu 0 führen. $\implies$ \textbf{Immer erst multiplizieren, dann dividieren!}
  \[
  \text{Falsch: } (a / b) * c \qquad \text{Richtig: } (a * c) / b
  \]
  \item \textbf{Quantisierungsfehler:} Bei kleinen Signalen führt Rundung zu Totzonen oder Grenzschwingungen.
\end{enumerate}
\end{tcolorbox}

\subsection{Anti-Windup-Algorithmen}
\begin{tcolorbox}
\boxtitle{Das Windup-Problem}
Bei Begrenzung der Stellgröße (z.\,B. Spannung $U_{\max}$, maximaler PWM-Wert) kann das System dem Sollwert nicht sofort folgen. Ein integrierender Regler (I-Anteil) summiert die Regeldifferenz weiter auf (\textit{Windup}). Beim Erreichen des Sollwerts ist der I-Anteil stark überladen und verursacht enormes Überschwingen.
\end{tcolorbox}

\begin{tcolorbox}
\boxtitle{1. Clamping (Integrator-Verriegelung)}
Das Aufintegrieren wird blockiert, wenn:
\begin{itemize}
  \item die Stellgröße in der Sättigung ist ($|u| \ge u_{\max}$) \textbf{und}
  \item die Regeldifferenz $e$ das gleiche Vorzeichen wie $u$ hat (was die Sättigung weiter vertiefen würde).
\end{itemize}
\end{tcolorbox}

\begin{tcolorbox}
\boxtitle{2. Back-Calculation (Rückrechnung)}
Die Differenz zwischen unbegrenzter Stellgröße $u$ und begrenztem Stellwert $u_{\text{sat}}$ wird über eine Nachführkonstante $\frac{1}{T_t}$ auf den Integrator rückgeführt:
\[
\dot{x}_I(t) = e(t) - \frac{1}{T_t} \cdot [u(t) - u_{\text{sat}}(t)]
\]
Im Zeitdiskreten (mit Abtastzeit $T_A$):
\[
x_{I,k} = x_{I,k-1} + T_A \cdot e_k - \frac{T_A}{T_t} \cdot (u_k - u_{\text{sat},k})
\]
Typische Wahl: $T_t \approx T_I$ (Nachstellzeit des I-Anteils).
\end{tcolorbox}

\subsection{Task-Scheduling \& Echtzeitanforderungen}
\begin{tcolorbox}
\boxtitle{Zustände eines Tasks}
\begin{itemize}
  \item \textbf{Running:} Task wird gerade von der CPU ausgeführt.
  \item \textbf{Ready:} Task ist ausführbereit und wartet auf CPU-Zuteilung durch den Scheduler.
  \item \textbf{Blocked / Waiting:} Task wartet auf ein externes Ereignis (Timer-Interrupt, Datenempfang).
\end{itemize}
\end{tcolorbox}

\begin{tcolorbox}
\boxtitle{Echtzeit-Anforderung für Regler}
\begin{itemize}
  \item Die Theorie zeitdiskreter Systeme setzt eine \textbf{strikt konstante Abtastzeit $T_A$} voraus.
  \item Zeitschwankungen (\textit{Jitter}) wirken wie additives Rauschen und Phasenverlust $\implies$ Destabilisierung!
  \item \textbf{Realisierung:} Der Regler-Task muss über einen \textbf{Hardware-Timer-Interrupt} mit \textbf{höchster Priorität} (preemptive) ausgeführt werden. Längere Hintergrundaufgaben (Display, Logging, Kommunikation) laufen mit niedrigerer Priorität im Hintergrund.
\end{itemize}
\end{tcolorbox}
